% Generated draft: requires mathematical and source review.
\documentclass[11pt]{article}
\usepackage[margin=25mm]{geometry}
\usepackage{fontspec}
\setmainfont{Latin Modern Roman}
\usepackage{amsmath,amssymb,mathtools,unicode-math}
\setmathfont{Latin Modern Math}
\usepackage{enumitem,needspace,xurl,hyperref}
\setlength{\parindent}{0pt}
\setlength{\parskip}{5pt}
\setlength{\emergencystretch}{3em}
\newcommand{\N}{\mathbb{N}}
\newcommand{\Z}{\mathbb{Z}}
\newcommand{\Q}{\mathbb{Q}}
\newcommand{\R}{\mathbb{R}}
\newcommand{\C}{\mathbb{C}}
\newcommand{\F}{\mathbb{F}}
\newcommand{\A}{\mathbb{A}}
\newcommand{\PP}{\mathbb{P}}
\newcommand{\E}{\mathbb{E}}
\newcommand{\eps}{\varepsilon}
\newcommand{\dd}{\,\mathrm{d}}
\begin{document}
\section*{UnsolvedMath ID: 10400107}
\textbf{Source problem code:} catalogue entry
\subsection{Definitions and mathematical statement}
A quandle \(Q\) has an idempotent operation \(*\), bijective right translations, and \((x*y)*z=(x*z)*(y*z)\). In the usual untwisted integral chain convention, let \(C_n^{\rm R}(Q)\) be the free abelian group on \(Q^n\). For \(n\ge2\), set
\[
 \partial_n(x_1,\ldots,x_n)=
 \sum_{i=2}^{n}(-1)^i\bigl(
 (x_1,\ldots,\widehat{x_i},\ldots,x_n)
 -(x_1*x_i,\ldots,x_{i-1}*x_i,x_{i+1},\ldots,x_n)\bigr).
\]
The degenerate subgroup is generated by tuples with \(x_i=x_{i+1}\) for some adjacent pair. Quotienting by these subgroups gives the quandle complex \(C_\bullet^{\rm Q}(Q)\). For an abelian coefficient group \(A\), ordinary untwisted quandle cohomology in this convention is
\[
 H^n_{\rm Q}(Q;A)=H^n\bigl(\operatorname{Hom}_{\Z}(C_\bullet^{\rm Q}(Q),A)\bigr).
\]
Degree-zero conventions and coefficient variants must be fixed separately; this definition is a conventional interpretation of the source term.

C. Rourke and B. Sanderson's Problem 5.11 asks whether there is a natural quandle space whose cohomology groups are the quandle cohomology groups. A usual precise interpretation would seek a space \(B_{\rm Q}(Q)\) assigned naturally to each quandle, together with isomorphisms
\[
 H^n(B_{\rm Q}(Q);A)\cong H^n_{\rm Q}(Q;A)
\]
in the intended degrees and coefficient convention, compatible with quandle homomorphisms. The source does not specify whether naturality is strict functoriality or only up to homotopy, which coefficients or degrees are included, or whether a particular cellular model, ring structure, or stronger chain-level realization is required. The displayed interpretation does not silently impose all these additional requirements or identify the requested space with a pre-existing rack-space construction.
\subsection{Short English statement}
Find a natural topological realization of quandle cohomology: a space associated to a quandle whose ordinary cohomology gives the intended quandle-cohomology groups. The source does not specify coefficients, low-degree conventions, or the required sense of naturality. Its supplied open-triage assessment awaits a current literature audit.
\subsection{Sources}
\begin{itemize}
\item Local UnsolvedMath registry, permanent ID \(10400107\), exact source code \texttt{catalogue entry}, Problem 5.11, attributed in the statement to C. Rourke and B. Sanderson. The title is truncated after ``(C.''
\item Ohtsuki, \emph{Problems on invariants of knots and 3-manifolds} (2002), Problem 5.11, PDF page 93: \url{https://msp.org/gtm/2002/04/gtm-2002-04-024s.pdf}. Bold printed headings distinguish 231 numbered items from surrounding references.
\item Extraction pdf-font-structure, accessed 2026-07-29. Source evidence NEEDS\_REVIEW presents an open problem and requires a current primary-literature audit. Supplied research classification OPEN-TRIAGE calls this natural-space question not obviously resolved. No further authoritative reference or proposed space is supplied.
\item Dataset attribution: Ulam AI / UnsolvedMath Contributors, CC BY 4.0. Source rights NEEDS\_REVIEW: publicly accessible author/publisher copy, redistribution terms requiring review.
\end{itemize}
\end{document}
